% Geometría de un alimentador de tornillo horizontal bajo una tolva.
% Elaboración propia con base en Bates (2000) y Karwat et al. (2021).
% Variables de la Ecuación (2.1): D, d, s y n.
\begin{tikzpicture}[
    >=Latex,
    font=\small,
    cota/.style={<->, thin},
    nota/.style={font=\footnotesize, align=center},
    flujo/.style={-{Latex[width=5pt,length=4pt]}, thick, gray}
]

% Tolva sobre la sección de alimentación
\draw[thick] (0.2,3.6) -- (1.0,1.6) (4.6,1.6) -- (5.4,3.6);
\node[nota] at (2.8,3.1) {Tolva};

% Artesa (canal) del tornillo
\draw[thick] (-0.6,1.6) -- (1.0,1.6) (4.6,1.6) -- (7.0,1.6);
\draw[thick] (-0.6,0.0) -- (7.0,0.0);
\draw[thick] (7.0,1.6) -- (7.0,1.0) (7.0,0.0) -- (7.0,0.3);

% Eje del tornillo (diámetro d)
\draw[thick, fill=gray!15] (-2.0,0.65) rectangle (7.4,0.95);

% Álabes (diámetro D), paso s entre álabes consecutivos
\foreach \x in {0.0,1.4,2.8,4.2,5.6}{
    \draw[very thick] (\x,0.1) -- (\x+0.7,1.5);
}

% Motor
\draw[thick, fill=gray!30] (-3.2,0.3) rectangle (-2.0,1.3);
\node[nota] at (-2.6,0.8) {Motor};
\draw[->, thick] (-1.75,1.15) arc[start angle=150, end angle=30, radius=0.3];
\node[nota] at (-1.5,1.65) {$n$};

% Descarga
\draw[flujo] (7.0,0.65) -- (7.0,-0.75);
\node[nota, anchor=west] at (7.1,-0.5) {$\dot{m}$};

% Cota del paso s
\draw[cota] (1.4,-0.35) -- (2.8,-0.35);
\node[nota] at (2.1,-0.6) {$s$};
\draw[thin, gray] (1.4,0.1) -- (1.4,-0.45) (2.8,0.1) -- (2.8,-0.45);

% Cota del diámetro exterior D
\draw[cota] (7.8,0.1) -- (7.8,1.5);
\node[nota, anchor=west] at (7.9,0.8) {$D$};
\draw[thin, gray] (6.3,1.5) -- (7.9,1.5) (6.3,0.1) -- (7.9,0.1);

% Cota del diámetro del eje d
\draw[cota] (-0.9,0.65) -- (-0.9,0.95);
\node[nota] at (-0.9,0.4) {$d$};

\end{tikzpicture}
